%%%%%%%%%%  scan idx 557  |  folio 544  %%%%%%%%%%
\SGAPrintPage{544}%

\sgaStableTarget{sga6.e10.sec.4.4}\subsection*{4.4.}
The integer (4.3.4) is sometimes called the \emph{intersection number} of the invertible \(\underline{O}_X\)-modules \(\underline{L}_1,\ldots,\underline{L}_d\) with \(\underline{F}\), and is denoted by \((\underline{L}_1,\ldots,\underline{L}_d;\underline{F})\) (cf. [3]). Expression (4.3.4) for this number immediately gives the following properties:

\begin{enumerate}[label=\arabic*)]
\item The integer (4.3.4) is a symmetric, \(d\)-linear function of the variables \(L_1,\ldots,L_d\in\Pic(X)\), zero when \(\dim(\Supp(\underline{F}))<d\), and is ``additive'' in \(\underline{F}\), i.e. additive in \(\cl_\bullet(\underline F)=F\).

\item Using 2.1.2, one immediately concludes (with the notation of 1.1.2):
\sgaStableTarget{sga6.e10.eq.4-4-1}\begin{equation}\tag{4.4.1}
 (\underline{L}_1,\ldots,\underline{L}_d;\underline{F})
 =
 \sum_{x\in X(d)}
 \operatorname{long}_{\underline{O}_{X,x}}(\underline{F}_x)
 (\underline{L}_1,\ldots,\underline{L}_d;\underline{O}_{Z(x)}).
\end{equation}

\item Let \(Y\) be a closed subscheme of \(X\) and \(\underline{G}\) a coherent \(\underline{O}_Y\)-module such that \(\underline{F}=f_*(\underline{G})\), where \(f:Y\to X\) is the canonical immersion. Then one obtains
\[
 (f^*\underline{L}_1,\ldots,f^*\underline{L}_d;\underline{G})=(\underline{L}_1,\ldots,\underline{L}_d;\underline{F}).
\]
More generally, the projection formula gives
\sgaStableTarget{sga6.e10.eq.4-4-2}\begin{equation}\tag{4.4.2}
 \chi(f^*(x)\cdot y)=\chi\bigl(f_*(f^*(x)\cdot y)\bigr)=\chi(x\cdot f_*(y))
 \qquad (x\in K^{\bullet}(X),\ y\in K_\bullet(Y))
\end{equation}
for every \(k\)-morphism \(f:Y\to X\) of proper \(k\)-schemes.

\item Suppose that, in the morphism \(f:Y\to X\) above, \(X\) and \(Y\) are irreducible of dimension \(\leq d\). If \(f\) has degree \(n\), then one has


%%%%%%%%%%  scan idx 558  |  folio 545  %%%%%%%%%%
\SGAPrintPage{545}%

\sgaStableTarget{sga6.e10.eq.4-4-3}\begin{equation}\tag{4.4.3}
 (f^*\underline{L}_1,\ldots,f^*\underline{L}_d;O_Y)
 =
 n(\underline{L}_1,\ldots,\underline{L}_d;\underline{O}_X).
\end{equation}
Recall the definition of the degree \(n\) of \(f\) [3]: it is zero except when \(\dim(X)=\dim(Y)=\dim(f(Y))\). In this latter case, put
\sgaStableTarget{sga6.e10.eq.4-4-4}\begin{equation}\tag{4.4.4}
 n=(\len_{\underline{O}_{X,x}}(\underline{O}_{Y,y}))/(\len_{\underline{O}_{X,x}}(\underline{O}_{X,x})),
\end{equation}
where \(x\) is the generic point of \(X\) and \(y\) the generic point of \(Y\).
\end{enumerate}

In the proof of (4.4.3), the case \(\dim(X)=\dim(Y)\) is the only nontrivial one. If \(\dim(f(Y))<d\), then one has
\sgaStableTarget{sga6.e10.eq.4-4-5}\begin{equation}\tag{4.4.5}
 R^if_*(\underline{O}_Y)=0
\end{equation}
for \(i>0\) in a neighborhood of \(x\). But condition (4.4.5) also holds when \(\dim(f(Y))=d\), since the generic fiber then has dimension zero (EGA III 4.2.2 and EGA IV 5.6.6). Then 1.1.2 gives the congruences
\[
 f_*(\cl_\bullet(\underline{O}_Y))
 \equiv \cl_\bullet(f_*(\underline{O}_Y)) \equiv
 \len_{\underline{O}_{X,x}}(\underline{O}_{Y,y})
 \cl_\bullet(\underline{O}_{X_{\mathrm{red}}})
\]
modulo \(\Filt_{d-1}(X)\). By (4.4.2), it follows that
\[
 (f^*\underline{L}_1,\ldots,f^*\underline{L}_d;\underline{O}_Y)
 =
 \len_{\underline{O}_{X,x}}(\underline{O}_{Y,y})
 (\underline{L}_1,\ldots,\underline{L}_d;\underline{O}_{X_{\mathrm{red}}}).
\]
On the other hand, \(2^\circ\) gives
\[
 (\underline{L}_1,\ldots,\underline{L}_d;\underline{O}_X)
 =
 \len_{\underline{O}_{X,x}}(\underline{O}_{X,x})
 (\underline{L}_1,\ldots,\underline{L}_d;\underline{O}_{X_{\mathrm{red}}}),
\]
whence the assertion.


%%%%%%%%%%  scan idx 559  |  folio 546  %%%%%%%%%%
\SGAPrintPage{546}%

\sgaStableTarget{sga6.e10.sec.4.5}\subsection*{4.5}
The equivalence relation defined by the intersection product \(\langle\, ,\,\rangle\) is called numerical equivalence:

\sgaStableTarget{sga6.e10.def.4-5-1}\textbf{Definition 4.5.1.} \uline{Two elements} \(x\) \uline{and} \(y\) \uline{of} \(\Gr^i(X)\) \((i\in Z)\) (resp. \uline{of} \(\Gr_i(X)\)\uline{) are said to be numerically equivalent when they define the same homomorphism} \(\Gr_i(X)\longrightarrow\mathbb Z\) (resp. \(\Gr^i(X)\longrightarrow\mathbb Z\)\uline{). We say that two invertible} \(\underline{O}_X\)\uline{-modules} \(\underline{L}_1\) and \(\underline{L}_2\) \uline{are numerically equivalent when their first Chern classes} \(x,y\in\Gr^1(X)\) \uline{are numerically equivalent.}

\sgaStableTarget{sga6.e10.prop.4-5-2}\textbf{Proposition 4.5.2.} \uline{An invertible} \(\underline{O}_X\)\uline{-module} \(\underline{L}\) \uline{is numerically equivalent to zero if and only if}
\sgaStableTarget{sga6.e10.eq.4-5-3}\begin{equation}\tag{4.5.3}
 \chi(\underline{L}^n\otimes_{\underline{O}_X}\underline{O}_Z)=\chi(\underline{O}_Z)
\end{equation}
\uline{for every} \(n\in\mathbb Z\) \uline{and every integral closed curve} \(Z\subset X\).

This follows from 1.1.2 and 4.3.1.

Denote by \(\Pic^N(X)\) the subgroup of \(\Pic(X)\) consisting of the classes of invertible \(\underline{O}_X\)-modules numerically equivalent to zero. From 2.4.4 one obtains

\sgaStableTarget{sga6.e10.cor.4-5-4}\textbf{Corollary 4.5.3.} \(\Pic^\tau(X)\subset \Pic^N(X)\).

In Exp. XIII it will be proved that one even has \(\Pic^N(X)=\Pic^\tau(X)\).

\sgaStableTarget{sga6.e10.sec.5}\section*{5. The isomorphism \(\Pic(X)\cong\Gr^1(X)\)}

\sgaStableTarget{sga6.e10.sec.5.1}\subsection*{5.1.}
Let \(X\) be an arbitrary ringed topos. Then one can define the ring \(K^{\bullet}(X)\) and the filtration \((\Filt^i(X))_{i\in\mathbb Z}\) exactly as for schemes. Likewise, define the group \(\Pic(X)\) as the group of isomorphism classes of invertible \(\underline{O}_X\)-modules. Denote by


%%%%%%%%%%  scan idx 560  |  folio 547  %%%%%%%%%%
\SGAPrintPage{547}%

\sgaStableTarget{sga6.e10.eq.5-1-1}\begin{equation}\tag{5.1.1}
 i:\Pic(X)\longrightarrow K^{\bullet}(X)
\end{equation}
the map associating to the class \(\cl_P(\underline{L})\in\Pic(X)\) of an invertible \(\underline{O}_X\)-module \(\underline{L}\) the element \(\cl^{\bullet}(\underline{L})\) of \(K^{\bullet}(X)\). This map is clearly a homomorphism from \(\Pic(X)\) to the multiplicative group \(K^{\bullet}(X)^*\) of invertible elements of \(K^{\bullet}(X)\). Henceforth, the group law of \(\Pic(X)\) will be written multiplicatively.

The map \(i\) has a left inverse defined by the ``det'' functor: when \(\underline{E}\) is a locally free \(\underline{O}_X\)-module of rank \(\underline{n}\), one has
\[
 \det(\underline{E})=\bigwedge^{\underline{n}} \underline{E}.
\]
If
\sgaStableTarget{sga6.e10.eq.5-1-2}\begin{equation}\tag{5.1.2}
0\longrightarrow \underline{E}'\longrightarrow \underline{E}\longrightarrow \underline{E}''\longrightarrow0
\end{equation}
is an exact sequence of locally free \(\underline{O}_X\)-modules, there is a canonical isomorphism
\[
 \det(\underline{E})\xrightarrow{\sim}(\det(\underline{E}'))\otimes_{\underline{O}_X}(\det(\underline{E}''))
\]
arising from (6.2.9) and the fact that (5.1.2) is locally split. Thus one obtains a homomorphism of abelian groups
\sgaStableTarget{sga6.e10.eq.5-1-3}\begin{equation}\tag{5.1.3}
 \det:K^{\bullet}(X)\longrightarrow\Pic(X)
\end{equation}
that is the required left inverse of \(i\). In particular, \(i\) is injective, and henceforth we shall identify \(\Pic(X)\) with its image in \(K^{\bullet}(X)^*\).

\sgaStableTarget{sga6.e10.sec.5.2}\subsection*{5.2.}
Define another map
\sgaStableTarget{sga6.e10.eq.5-2-1}\begin{equation}\tag{5.2.1}
 \delta:\Pic(X)\longrightarrow\Filt^1(X)
\end{equation}


%%%%%%%%%%  scan idx 561  |  folio 548  %%%%%%%%%%
\SGAPrintPage{548}%

by putting, for every \(L\in\Pic(X)\),
\[
 \delta(L)=1-L^{-1}.
\]

\sgaStableTarget{sga6.e10.rem.5-2-2}\textbf{Remarks 5.2.2.}:
\begin{enumerate}[label=\alph*)]
\item Suppose that \(X\) is the locally ringed topos of Zariski sheaves on a scheme \(S\). When \(D\) is a Cartier divisor \(\geq0\) on \(X\) and
\[
 L=\cl^{\bullet}(\underline{O}_S(D)),
\]
one has
\sgaStableTarget{sga6.e10.eq.5-2-3}\begin{equation}\tag{5.2.3}
 \theta_S(\delta(L))=\cl_\bullet(\underline{O}_D)
\end{equation}
where \(\theta_S:K^{\bullet}(X)\to K_\bullet(S)\) is the canonical homomorphism of 1.2.3. This follows immediately from the exact sequence
\[
0\longrightarrow\underline{O}_S(-D)\longrightarrow\underline{O}_S\longrightarrow\underline{O}_D\longrightarrow0.
\]

\item Under the hypotheses of 5.1 and 5.2, one has, for every \(L\in\Pic(X)\),
\sgaStableTarget{sga6.e10.eq.5-2-4}\begin{equation}\tag{5.2.4}
 \delta(L)\equiv L-1\quad\operatorname{mod.}\Filt^2(X),
\end{equation}
which justifies the notation \(c_1(L)\) introduced below. Indeed, one has
\[
 (L-1)-(1-L^{-1})=(L-1)(1-L^{-1})\in\Filt^2(X).
\]
\end{enumerate}

\sgaStableTarget{sga6.e10.sec.5.3}\subsection*{5.3.}
Composing \(\delta\) with the canonical epimorphism
\[
 \Filt^1(X)\longrightarrow\Gr^1(X)
\]
gives the ``first Chern class'' on \(\Pic(X)\):
\sgaStableTarget{sga6.e10.eq.5-3-1}\begin{equation}\tag{5.3.1}
 c_1:\Pic(X)\longrightarrow\Gr^1(X).
\end{equation}

